\section{Manuscript corrections and integration}
\label{sec:corrections}

\subsection{Replace the broad monotonicity suggestion by a quantified question}

The final research chapter still says that monotonicity of higher zero
branches in the deformation parameter is suggested by the first
correction \cite{ProjectDiscovery}. Read literally as a direction for
all indices, that suggestion is obsolete once the incoming results are
taken into account. The earlier boundary and bifurcation reports already
give a turning lower branch at $(n,k)=(2,1)$, while the global-phase
report proves an increasing middle branch at $(3,1)$
\cite{IncomingBoundary,IncomingBifurcation,IncomingPhase}.
The positive initial velocity at $(n,k)=(4,3)$ also appears in the
formal-reductions report \cite{IncomingReductions}.

The additional correction supplied here is infinite in scope:
Theorem~\ref{diag:thm:oscillation} gives both directions at arbitrarily
large indices on $k=n-1$. A replacement paragraph should therefore
ask for \emph{effective derivative-order thresholds depending on the
Laurent index}, the distribution of turning points below those
thresholds, and the joint-index scales on which a uniform direction
does emerge. It should not suggest that every higher branch decreases
throughout the deformation.

One precise consequence is Corollary~\ref{diag:cor:threshold}:
if $K_n$ guarantees decreasing motion for every branch and every
$k\ge K_n$, then $K_n>n-1$ for infinitely many $n$.
This does not disprove the existence of $K_n$, does not show that
$K_n/n$ is unbounded, and does not identify a saturation threshold
for the number of zeros. Those are different claims.

\subsection{Record the exact endpoint in the optimized barrier}

The solution of the incoming exterior-operator problem should be recorded
in two parts. The least uniform cutoff for the summed family is attained
and its inequality is strict when $\rho>0$. The strict pointwise cutoff
has the same infimum, but the infimum is not attained. At the optimal
coefficient,
\[
 g_n'(x)+\lambda_ng_n(x)=0
 \quad\Longleftrightarrow\quad
 x\in\{B_n,e^{n+r_n}\},\qquad x\ge B_n.
\]
Consequently the strict summed theorem must not be extended unchanged to
$\rho=0$, to arbitrary one- or two-point positive measures, or to a
claim that every individual summand is strictly negative.
Theorem~\ref{opt:thm:optimal} and its positive-measure extension state
the exact hypotheses. This is a qualification of the newly optimized
result, not an assertion that the earlier fixed-coefficient theorem
was false.

\subsection{Keep integer special values separate from order-zero geometry}

For reference, write
$g_{a,b}=\operatorname{Im}\operatorname{Li}_{a,b}(i,1)$ and
$G=\beta(2)$. The remaining manuscript conjecture is
\begin{align}
 S(6)\stackrel{?}{={}}&
 \frac{722}{527}g_{6,1}+\frac{40}{527}g_{4,3}
 -\frac{128}{155}g_{2,5}+\frac{15191}{28569600}\pi^7
 \notag\\
 &-\frac3{124}G\zeta(5)
 -\frac{2373}{10540}\beta(4)\zeta(3)
 -2\beta(6)\log2 .
 \label{eq:retained-S6}
\end{align}
This is quoted to identify the open target, not proposed as a new
identity. The exact shuffle relation already in the manuscript is
\begin{align}
 g_{2,5}={}&30g_{6,1}+15g_{5,2}+5g_{4,3}
 -\frac{15}{512}G\zeta(5)\notag\\
 &+\frac3{16}\beta(4)\zeta(3)-\frac{11\pi^7}{368640}.
 \label{eq:retained-basket-change}
\end{align}
It proves equivalence of the two reported $S_6$ baskets
\cite{ProjectQuotients}. It does not prove their residual vanishes.

Our Fourier--beta identity gives new exact representations of $S(6)$
by analytic continuation, but does not reduce it to the specified
weight-seven period basket. Entirety and a complete real zero set
likewise do not imply that reduction. The finite separating functional
in the manuscript obstructs a particular restricted
shuffle/stuffle presentation; it does not prove arithmetic independence.
The successful $S_4$ certificate shows why a convergent octahedral
relation and suitable higher-depth lifts are natural next ingredients
\cite{ProjectS4,ProjectQuotients}.

The incoming rejected $S_8$ vector is a useful methodological boundary:
an interval containing zero is compatible with equality, but does not
prove it; a certified interval excluding zero rejects the specific
frozen vector, but does not prove independence of its basket.
The proposed integration preserves those distinctions
\cite{ProjectDiscovery}.

\subsection{Notation and reusable proof components}

The generalized secant numbers in
Corollary~\ref{sorder:cor:negative-values} are positive at $\alpha=1$.
They differ by $(-1)^m$ from the signed Euler numbers sometimes defined
through $\sech t$. Their generating function, not the letter $E$ alone,
fixes the convention.
The symbols $S_p$ and $S(p)$ agree at positive integers. At negative
integers $S(s)$ always denotes analytic continuation, since the
original alternating Dirichlet series no longer converges.

The harmonic continuation and its Fourier identity can be integrated
after the Gaussian harmonic sums. The diagonal theorem belongs after
the elementary endpoint and fixed-index Lerch zero geometry. The
optimized barrier should follow the incoming global-phase theorem,
with its research question~2 marked resolved with the attainment
qualification above. A suggested replacement paragraph and dependency
map are included in the integration artifacts.

\section{A proposed research programme}
\label{sec:programme}

The results suggest several specific next problems. Each item below
states what is missing and what would count as progress.

\subsection{The order interpolation and its zeros}

\paragraph{1. Nonreal zeros and growth of the entire function.}
Determine the growth order and directional indicator of $S(s)$, then
locate and count its nonreal zeros. The real classification and the
zero-free half-plane $\Re s\ge3$ are starting constraints, not a
classification in the complex plane. A useful first target is a rigorous
argument-principle count in bounded rectangles using regularized Mellin
integrals and controlled tails. A global conclusion would require
uniform complex asymptotics.

\paragraph{2. Explicit finite-index errors for $\sigma_m$.}
Replace the implicit constants in
Theorem~\ref{sorder:thm:zero-asymptotic} by computable bounds valid for
every $m$ beyond a stated threshold. The proof already reduces this
to a sectorial digamma remainder and beta-series estimates.
Combining those with a lower bound on the derivative of the tangent
comparison would give certified intervals of predictable width
without numerical root scanning.

\paragraph{3. Shifted denominators and higher harmonic depth.}
Study
\[
 S_a(s)=\sum_{n\ge0}\frac{(-1)^nH_n}{(n+a)^s},
 \qquad a>0,
\]
and the elementary-symmetric harmonic sums already used by the
manuscript. Their Mellin kernels are explicit, but the self-reflecting
choice $a=1/2$ is special: its Fourier transform gives the positive
moment ratio used here. Determine which shifts preserve a comparable
sign or total-positivity structure. It would be premature to assert
the same real zero pattern for every shift.

\paragraph{4. Sharper moment-ratio inequalities.}
The estimate $R'(q)<\sqrt{181/27}/q$ is intentionally coarse.
Find the best constant or a useful pointwise comparison, and determine
whether $R'$ decreases on $[2,\infty)$. The proof establishes neither
concavity nor complete monotonicity of $R$. A refined covariance
inequality may be more efficient than additional moment bounds.

\subsection{Joint-index Lerch geometry}

\paragraph{5. Prove or refute the accessible-singularity conjecture.}
The exact inverse equation for $u_A$ makes its finite critical values
explicit, but those values belong to different analytic sheets.
For $A=2$, prove which critical points are accessible from the germ
at zero, identify all singularities on its convergence circle, and
establish the continuation domain needed for coefficient transfer.
This would turn the quantitative conjecture into an asymptotic theorem.
Numerical radial continuation and agreement of many coefficients
are evidence, not a substitute for that domain argument.

\paragraph{6. Density and gaps of the two velocity signs.}
Theorem~\ref{diag:thm:oscillation} proves infinitude of both signs without
a density or a maximum gap. A dominant-conjugate-pair theorem should
give an oscillatory law, but a density statement also depends on the
arithmetic and possible resonances of its phase. For algebraic
$A>1$, the finite Stirling formula already excludes a zero velocity;
it does not bound how small a nonzero velocity may be.

\paragraph{7. Effective monotonicity thresholds.}
Determine the growth of the least $K_n$ for which all zero branches
decrease for $k\ge K_n$. The present lower constraint
$K_n>n-1$ along infinitely many indices should be combined with the
existing fixed-$n$ concentration estimates. A linear upper bound,
a sharp transition ratio, or a proof that no linear bound is possible
would each be meaningful outcomes; none follows merely from the
diagonal sign theorem.

\paragraph{8. Central root splitting for $n=k+r$.}
For fixed $r>1$, the elementary root at $a=1$ has multiplicity $r$.
Its natural perturbation can involve fractional powers of $\rho$.
Derive a uniform local normal form, determine the number of real
branches it produces, and compare it with the existing expansions
near positive Bessel zeros \cite{IncomingBessel}. The central root
and positive Bessel roots should not be treated as interchangeable
regimes.

\paragraph{9. Beyond first-order velocity.}
Find generating functions for higher $\rho$ derivatives of
$a_n^{(h)}(\rho)$ and for the size of the neighborhood in which
the first derivative determines its motion. Small nonzero
$v_n$ can force second-order terms to become relevant unusually early.
This is a concrete way to connect the local oscillation theorem
to actual turning points in the full deformation.

\subsection{Exterior barriers and remaining phase questions}

\paragraph{10. Fixed-parameter and variable-coefficient barriers.}
Theorem~\ref{opt:thm:optimal} optimizes a constant coefficient uniformly
for arbitrarily small positive $\rho$. For a fixed positive $\rho$, the tail sum may allow
a smaller cutoff. Determine the optimum as a function of $\rho$,
or allow a controlled positive function $\lambda(a)$.
Such an extension must specify its admissible class; otherwise a
best-operator question is not a well-defined optimization problem.

\paragraph{11. Existence inside the exterior uniqueness region.}
The optimal barrier proves at most one exterior zero and a downward
crossing. It does not guarantee a zero. Prove signs at $B_n$ for useful
ranges of $(n,\rho)$, perhaps using a small set of exact separators
and geometric-weight inequalities. This would turn a one-crossing
bound into an all-index existence and uniqueness theorem.

\paragraph{12. Strict span unimodality at index three.}
The incoming report already proves the exact unit maximum and its
nondegeneracy, but its stronger conjecture of strict increase before
the maximum and strict decrease afterwards remains open
\cite{IncomingPhase}. Additional control of the averaged outer slope
is required; the sign of the shifted-root identity alone does not
settle this shape question.

\subsection{Arithmetic reductions and formal verification}

\paragraph{13. A weight-seven mixed certificate.}
For \eqref{eq:retained-S6}, seek a convergent octahedral relation whose
image is nonzero under the known restricted separating functional.
Then lift and eliminate depth-three or higher terms with exact rational
rows. The useful output is a short certificate with an independently
specified word convention and regularization, rather than another
decimal relation fit. A failed search in a finite row space should be
recorded as a limitation of that row space.

\paragraph{14. Formalize the finite and analytic interfaces separately.}
The Stirling-polynomial recurrence, cubic factorization, rational
interval operations, and negative-integer recurrences are finite
targets suited to proof-assistant formalization. The subsequent
analytic targets are Mellin subtraction, the one-dimensional variance
inequality, the positive-coefficient argument, and the implicit
zero comparison. The current package provides ordinary proofs and
executable checks; it does not claim a completed Lean or other
proof-assistant verification.
